\section{Các công trình tiêu biểu trong tutorial}
\label{sec:related-papers}

Các chương trước đã trình bày nền tảng, các nhóm jailbreak attacks,
jailbreak defenses, phương pháp đánh giá và sự mở rộng của bài toán sang
agentic systems. Chương này tổng hợp và phân tích các công trình tiêu biểu
được đề cập trong tutorial ICML 2025 \textit{Jailbreaking LLMs and Agentic
Systems: Attacks, Defenses, and Evaluations}. Mục tiêu của chương không
phải là bao quát toàn bộ literature về jailbreaking, mà là làm rõ vai trò
của các paper trung tâm trong từng nhóm nội dung: attack, defense,
evaluation và agentic/robotic systems.

Khác với Chương~\ref{sec:jailbreak-attacks} và
Chương~\ref{sec:jailbreak-defenses}, chương này không đi sâu vào cơ chế kỹ
thuật của từng nhóm phương pháp. Thay vào đó, mỗi công trình được đặt vào
bối cảnh chung của tutorial để làm rõ câu hỏi: paper đó đại diện cho hướng
nghiên cứu nào, đóng góp chính là gì, và tại sao nó quan trọng đối với bài
toán jailbreak.

\subsection{Các công trình về jailbreak attacks}
\label{subsec:attack-papers}

Nhóm attack papers trong tutorial cho thấy sự phát triển của jailbreak từ
các prompt thủ công sang các phương pháp tự động, có khả năng tối ưu, tìm
kiếm, chuyển giao và thích nghi với mô hình mục tiêu.

\subsubsection{Universal and Transferable Adversarial Attacks}

Công trình \textit{Universal and Transferable Adversarial Attacks on
Aligned Language Models} là một trong những paper trung tâm nhất của phần
attack \cite{zou2023universal}. Tutorial sử dụng công trình này để minh
họa hướng white-box jailbreak dựa trên adversarial suffix và thuật toán
Greedy Coordinate Gradient (GCG).

Ý tưởng chính của paper là tìm một chuỗi token bổ sung được nối vào yêu
cầu ban đầu nhằm làm suy yếu hành vi từ chối của mô hình. Nếu $x$ là yêu
cầu gốc và $a$ là adversarial suffix, đầu vào sau tấn công có dạng

\begin{equation}
    x' = x \oplus a,
    \label{eq:chapter7-gcg-input}
\end{equation}

trong đó $\oplus$ là phép nối chuỗi. Suffix $a$ được tối ưu để làm tăng
khả năng mô hình tạo phản hồi theo hướng attacker mong muốn. Về mặt khái
quát, bài toán có thể được biểu diễn bởi

\begin{equation}
    a^{*}
    =
    \arg\min_{a \in \mathcal{V}^{m}}
    \mathcal{L}
    \left(
        F_{\theta}(x \oplus a)
    \right),
    \label{eq:chapter7-gcg-objective}
\end{equation}

trong đó $\mathcal{V}$ là vocabulary, $m$ là độ dài suffix và
$\mathcal{L}$ là hàm mất mát dùng để hướng đầu ra của mô hình.

Vai trò của công trình này trong tutorial là làm rõ ba điểm. Thứ nhất,
jailbreak có thể được xây dựng bằng tối ưu hóa, không chỉ bằng prompt thủ
công. Thứ hai, các chuỗi adversarial có thể có hình thức không tự nhiên
đối với con người nhưng vẫn tạo tác động lớn lên mô hình. Thứ ba, một số
attack có khả năng transfer giữa các mô hình, cho thấy các mô hình khác
nhau có thể chia sẻ điểm yếu trong cơ chế safety alignment.

\subsubsection{PAIR}

Công trình \textit{Jailbreaking Black Box Large Language Models in Twenty
Queries} giới thiệu PAIR, một phương pháp black-box automated jailbreak
\cite{chao2023pair}. Đây cũng là một paper được tutorial trình bày sâu vì
nó minh họa cách jailbreak có thể được tự động hóa ngay cả khi attacker
không có quyền truy cập tham số, gradient hoặc logit của mô hình.

PAIR sử dụng ba thành phần chính: attacker model, target model và judge.
Attacker model sinh candidate prompt; target model phản hồi; judge đánh
giá mức độ thành công; sau đó attacker model tinh chỉnh prompt dựa trên
kết quả quan sát được. Quy trình này có thể được mô tả khái quát bởi

\begin{equation}
    x_{t+1}
    =
    A_{\phi}
    \left(
        x_t,
        y_t,
        J(x_t, y_t)
    \right),
    \label{eq:chapter7-pair-update}
\end{equation}

trong đó $A_{\phi}$ là attacker model, $x_t$ là prompt tại vòng lặp thứ
$t$, $y_t$ là phản hồi của target model và $J$ là evaluator.

Đóng góp chính của PAIR là chuyển jailbreak từ một hoạt động thủ công sang
một vòng lặp automated red teaming. Paper này cũng làm rõ một threat model
quan trọng trong thực tế: nhiều hệ thống thương mại chỉ cho phép truy cập
black-box qua API, nhưng điều đó không ngăn attacker sử dụng một mô hình
khác để tự động sinh và cải tiến prompt tấn công.

\subsubsection{Tree of Attacks}

\textit{Tree of Attacks: Jailbreaking Black-Box LLMs Automatically} mở
rộng hướng black-box automated jailbreak bằng cách sử dụng tìm kiếm dạng
cây \cite{mehrotra2023tap}. Nếu PAIR tinh chỉnh prompt theo một chuỗi lặp,
thì Tree of Attacks duy trì nhiều candidate prompts song song, mở rộng các
nhánh có triển vọng và loại bỏ các nhánh kém hiệu quả.

Có thể mô tả tập candidate ở bước $t$ là

\begin{equation}
    \mathcal{C}_t
    =
    \{x_t^{(1)}, x_t^{(2)}, \ldots, x_t^{(k)}\}.
    \label{eq:chapter7-tap-candidates}
\end{equation}

Sau bước mở rộng và pruning, hệ thống giữ lại tập con có khả năng thành
công cao hơn:

\begin{equation}
    \mathcal{C}_{t+1}
    =
    \operatorname{Prune}
    \left(
        \operatorname{Expand}(\mathcal{C}_t)
    \right).
    \label{eq:chapter7-tap-expand-prune}
\end{equation}

Vai trò của paper này trong tutorial là minh họa rằng black-box jailbreak
không chỉ có thể được tự động hóa, mà còn có thể được tổ chức như một bài
toán search. Điều này đặc biệt quan trọng khi ngân sách truy vấn bị giới
hạn: một chiến lược tìm kiếm tốt có thể giảm số lần thử kém hiệu quả và
tập trung vào các hướng có khả năng thành công cao hơn.

\subsubsection{AutoDAN}

\textit{AutoDAN: Generating Stealthy Jailbreak Prompts on Aligned Large
Language Models} đại diện cho hướng tự động tạo jailbreak prompt có hình
thức tự nhiên và khó phát hiện hơn so với adversarial suffix thông thường
\cite{liu2023autodan}. Nếu GCG có thể sinh ra các chuỗi token bất thường,
AutoDAN nhấn mạnh tính \textit{stealthy}: prompt được tạo ra cần vừa có khả
năng tấn công, vừa ít bị phát hiện bởi các defense dựa trên bất thường bề
mặt.

Vai trò của AutoDAN trong tutorial là bổ sung một góc nhìn quan trọng:
không phải mọi jailbreak tự động đều tạo ra chuỗi token khó đọc. Một attack
có thể được tối ưu để trông tự nhiên hơn, khiến input filtering hoặc các
heuristic phát hiện prompt lạ trở nên kém hiệu quả. Điều này làm rõ giới
hạn của các defense chỉ dựa trên đặc trưng bề mặt.

\subsubsection{Automated red teaming và reasoning-based search}

Tutorial cũng đề cập đến các hướng automated red teaming rộng hơn, chẳng
hạn \textit{Red Teaming Language Models with Language Models}
\cite{perez2022redteaming} và \textit{Adversarial Reasoning at Scale}. Các
công trình này đặt nền tảng cho ý tưởng dùng mô hình ngôn ngữ hoặc
test-time computation để tự động tìm các hành vi lỗi của một hệ thống khác.

Điểm chung của nhóm này là thay vì chỉ kiểm thử mô hình bằng một tập prompt
cố định, quá trình đánh giá được mở rộng thành một vòng lặp tìm kiếm chủ
động. Một mô hình hoặc một hệ thống search đóng vai trò red teamer, sinh ra
các tình huống kiểm thử, quan sát phản hồi của target model và tiếp tục mở
rộng các trường hợp có rủi ro. Đây là tiền đề quan trọng cho các phương
pháp như PAIR, TAP và các framework automated red teaming sau này.

\subsection{Các công trình về jailbreak defenses}
\label{subsec:defense-papers}

Nhóm defense papers trong tutorial cho thấy các phương pháp phòng thủ đã
phát triển từ input/output filtering đơn giản sang inference-time defense,
representation-level defense và các cơ chế can thiệp sâu hơn vào hành vi
mô hình.

\subsubsection{SmoothLLM}

\textit{SmoothLLM: Defending Large Language Models Against Jailbreaking
Attacks} là công trình trung tâm trong phần defense của tutorial
\cite{robey2023smoothllm}. SmoothLLM là một inference-time defense, tức là
nó không yêu cầu huấn luyện lại mô hình chính. Thay vào đó, phương pháp tạo
nhiều biến thể nhỏ của cùng một prompt và kiểm tra độ ổn định của phản hồi.

Với prompt đầu vào $x$, SmoothLLM tạo ra $K$ biến thể

\begin{equation}
    \tilde{x}_1, \tilde{x}_2, \ldots, \tilde{x}_K
    \sim
    \mathcal{P}(x),
    \label{eq:chapter7-smoothllm-perturbation}
\end{equation}

trong đó $\mathcal{P}$ là phân phối perturbation. Sau đó, hệ thống truy vấn
mô hình trên các biến thể này và tổng hợp kết quả:

\begin{equation}
    d
    =
    g
    \left(
        F_{\theta}(\tilde{x}_1),
        F_{\theta}(\tilde{x}_2),
        \ldots,
        F_{\theta}(\tilde{x}_K)
    \right).
    \label{eq:chapter7-smoothllm-aggregation}
\end{equation}

Ý tưởng cốt lõi là nhiều adversarial prompts, đặc biệt là các prompt được
tối ưu ở mức token hoặc ký tự, có tính giòn trước các biến đổi nhỏ.
SmoothLLM khai thác đặc điểm này để phát hiện hoặc làm giảm hiệu quả của
jailbreak.

Vai trò của SmoothLLM trong tutorial là minh họa một defense có thể triển
khai ở runtime, độc lập tương đối với mô hình chính. Tuy nhiên, tutorial
cũng nhấn mạnh hạn chế của hướng này: chi phí suy luận tăng do cần nhiều
lần gọi mô hình, và defense có thể bị thách thức bởi adaptive attacks hoặc
các prompt có tính ổn định ngữ nghĩa cao hơn.

\subsubsection{Circuit Breakers}

\textit{Circuit Breakers: Enhanced Safety in LLMs via Representation
Rerouting} đại diện cho nhóm model-level defense \cite{zou2024circuit}.
Thay vì chỉ kiểm tra input hoặc output, hướng này tìm cách can thiệp vào
biểu diễn bên trong của mô hình để ngăn mô hình đi vào các trạng thái biểu
diễn liên quan đến hành vi không an toàn.

Khác với SmoothLLM, Circuit Breakers không chỉ là một wrapper bên ngoài.
Nó hướng đến việc thay đổi cách mô hình xử lý các yêu cầu rủi ro ở tầng
representation. Về mặt trực giác, khi mô hình bắt đầu đi theo một hướng
kích hoạt có khả năng dẫn đến phản hồi không an toàn, defense sẽ reroute
hoặc can thiệp để đưa mô hình về hành vi an toàn hơn.

Vai trò của paper này trong tutorial là cho thấy defense có thể được đặt ở
mức sâu hơn so với input/output filtering. Tuy nhiên, các phương pháp dạng
này thường yêu cầu quyền truy cập vào mô hình hoặc quá trình huấn luyện,
nên khó áp dụng trực tiếp cho các API đóng.

\subsubsection{Representation Guardrails}

Nhóm \textit{Representation Guardrails} cũng thuộc hướng representation-
level defense. Thay vì chỉ dựa vào nội dung văn bản của prompt hoặc
response, các phương pháp này khai thác hidden representations, activation
patterns hoặc embedding space để phát hiện và kiểm soát hành vi không an
toàn.

Vai trò của nhóm này trong tutorial là bổ sung một tầng phòng thủ khác:
nếu attacker có thể làm prompt trông tự nhiên ở bề mặt, hệ thống vẫn có thể
phát hiện tín hiệu rủi ro trong biểu diễn bên trong của mô hình. Tuy nhiên,
giống Circuit Breakers, representation-level guardrails thường khó triển
khai hơn trong bối cảnh mô hình đóng và cần đánh giá cẩn thận để tránh làm
giảm utility.

\subsubsection{Instruction-following và utility evaluation}

Tutorial cũng nhắc đến hướng đánh giá ảnh hưởng của defense đến khả năng
làm theo chỉ dẫn hợp lệ. Đây là điểm quan trọng vì một defense không thể
chỉ được đánh giá bằng việc giảm Attack Success Rate. Nếu defense khiến mô
hình từ chối quá nhiều yêu cầu hợp lệ hoặc làm giảm khả năng instruction
following, hệ thống sẽ mất utility.

Do đó, các paper và benchmark về instruction-following evaluation được
dùng trong tutorial như một lời nhắc rằng defense phải được đánh giá theo
cả hai chiều: giảm unsafe compliance và hạn chế over-refusal.

\subsection{Các benchmark và phương pháp đánh giá}
\label{subsec:evaluation-papers}

Nhóm benchmark và evaluation papers trong tutorial giải quyết vấn đề: làm
thế nào để so sánh attack và defense một cách công bằng. Nếu mỗi công
trình dùng một tập prompt, target model, system prompt và evaluator khác
nhau, kết quả sẽ khó đối chiếu. Vì vậy, các benchmark như AdvBench,
HarmBench và JailbreakBench đóng vai trò quan trọng trong việc chuẩn hóa
đánh giá.

\subsubsection{AdvBench}

AdvBench là tập harmful behaviors được sử dụng phổ biến trong nghiên cứu
jailbreak, đặc biệt trong các công trình về adversarial suffix như GCG
\cite{zou2023universal}. Benchmark này cung cấp các yêu cầu rủi ro dạng
text-only để kiểm tra liệu attack có khiến mô hình tạo phản hồi không an
toàn hay không.

Vai trò của AdvBench trong tutorial là làm ví dụ cho một tập hành vi chuẩn
dùng để đánh giá attack. Tuy nhiên, AdvBench không bao phủ toàn bộ không
gian rủi ro, đặc biệt là multi-turn jailbreak, prompt injection trong
agentic systems hoặc tool-use attacks. Vì vậy, nó phù hợp cho đánh giá
ban đầu nhưng cần kết hợp với các benchmark rộng hơn.

\subsubsection{HarmBench}

\textit{HarmBench: A Standardized Evaluation Framework for Automated Red
Teaming and Robust Refusal} là benchmark quan trọng trong phần evaluation
\cite{mazeika2024harmbench}. HarmBench không chỉ là một tập prompt, mà là
một framework đánh giá có cấu trúc cho automated red teaming và robust
refusal.

HarmBench chuẩn hóa nhiều thành phần: behavior set, attack methods, target
models, defenses và evaluator. Điều này giúp kết quả giữa các phương pháp
dễ so sánh hơn. Trong tutorial, HarmBench đại diện cho nỗ lực đưa
jailbreak evaluation từ các thử nghiệm rời rạc sang quy trình đánh giá có
kiểm soát.

Điểm quan trọng của HarmBench là khái niệm \textit{robust refusal}. Một mô
hình không chỉ cần từ chối yêu cầu không an toàn trong điều kiện thông
thường, mà còn phải duy trì từ chối trước các biến thể attack khác nhau.

\subsubsection{JailbreakBench}

\textit{JailbreakBench: An Open Robustness Benchmark for Jailbreaking Large
Language Models} nhấn mạnh tính mở, khả năng tái lập và chuẩn hóa
benchmark \cite{chao2024jailbreakbench}. Benchmark này cung cấp attack
artifacts, behavior set, target models, scoring functions và leaderboard để
so sánh attack--defense trong một thiết lập rõ ràng.

Vai trò của JailbreakBench trong tutorial là làm rõ rằng evaluation không
chỉ là tính ASR trên một tập prompt. Một benchmark tốt cần công bố threat
model, system prompt, chat template, scoring rule và điều kiện đánh giá.
Điều này đặc biệt quan trọng trong nghiên cứu jailbreak vì kết quả có thể
thay đổi mạnh theo target model, system prompt, evaluator và ngân sách
truy vấn.

\subsubsection{Trojan Detection Challenge}

Trojan Detection Challenge (TDC) được tutorial nhắc đến như một nguồn hoặc
benchmark liên quan đến đánh giá hành vi tấn công, backdoor hoặc trojan.
TDC không phải trung tâm của jailbreak text-only, nhưng có giá trị như một
điểm liên hệ với các hướng an toàn mô hình rộng hơn.

Trong báo cáo này, TDC được xem là benchmark bổ trợ, giúp phân biệt
jailbreak tại thời điểm suy luận với các rủi ro khác như backdoor hoặc
trojan được cài trong mô hình. Nội dung này cũng nhấn mạnh rằng jailbreak
chỉ là một phần của bức tranh an toàn AI rộng hơn.

\subsection{Các công trình về robotics và agentic systems}
\label{subsec:agentic-robotics-papers}

Phần cuối của tutorial mở rộng bài toán từ chatbot sang agentic systems và
robotics. Trong bối cảnh này, rủi ro không chỉ nằm ở phản hồi văn bản mà
còn nằm ở kế hoạch, lời gọi công cụ, memory, observation và hành động trong
môi trường.

\subsubsection{RoboPAIR}

\textit{RoboPAIR: Jailbreaking LLM-Controlled Robots} là paper trung tâm
trong phần robotics của tutorial \cite{robey2024robopair}. Công trình này
mở rộng tinh thần của PAIR sang các hệ thống robot do LLM điều khiển. Thay
vì chỉ đánh giá xem mô hình có tạo phản hồi không an toàn hay không,
RoboPAIR xem xét liệu hệ thống robot có thể bị dẫn đến hành động trái với
ràng buộc an toàn hay không.

Vai trò của RoboPAIR trong tutorial là rất quan trọng vì nó thay đổi đơn
vị đánh giá của jailbreak. Với chatbot, output chính là văn bản. Với robot,
output có thể là hành động vật lý. Vì vậy, một defense chỉ kiểm tra phản
hồi cuối cùng là không đủ. Hệ thống cần kiểm soát kế hoạch, lời gọi hành
động, controller, môi trường và cơ chế dừng an toàn.

RoboPAIR cho thấy khi LLM được tích hợp vào embodied agents, jailbreak
không còn là vấn đề của ngôn ngữ đơn thuần. Nó trở thành vấn đề an toàn hệ
thống, trong đó controller, permission, sandbox, monitoring và human
confirmation đều đóng vai trò quan trọng.

\subsubsection{AgentDojo}

\textit{AgentDojo: A Dynamic Environment to Evaluate Prompt Injection
Attacks and Defenses for LLM Agents} là benchmark được tutorial nhắc đến
trong phần agentic systems \cite{debenedetti2024agentdojo}. AgentDojo tập
trung vào đánh giá prompt injection attacks và defenses trong môi trường
agent có khả năng sử dụng công cụ.

Điểm quan trọng của AgentDojo là nó đánh giá đồng thời utility và security.
Một agent phải hoàn thành nhiệm vụ hợp lệ của người dùng, nhưng không được
bị thao túng bởi dữ liệu hoặc chỉ dẫn không đáng tin cậy. Điều này phản
ánh đúng thách thức của agentic systems: defense quá yếu làm tăng security
violation, nhưng defense quá mạnh có thể làm agent không hoàn thành nhiệm
vụ.

Trong tutorial, AgentDojo đại diện cho hướng benchmark động, nơi agent
tương tác với môi trường và công cụ thay vì chỉ trả lời một prompt tĩnh.

\subsubsection{BrowserART và Breaking ReAct}

BrowserART và \textit{Breaking ReAct} được tutorial dẫn như các công trình
liên quan đến bảo mật của web/browser agents và ReAct agents. Các hệ thống
này đặc biệt dễ gặp rủi ro vì agent phải quan sát nội dung từ môi trường
bên ngoài, suy luận dựa trên observation và gọi công cụ hoặc hành động tiếp
theo.

\textit{Breaking ReAct} làm rõ một điểm quan trọng: trong các agent dùng mô
hình observation--reasoning--action, attacker có thể tác động không chỉ vào
prompt ban đầu mà còn vào observation hoặc tool output. Nếu agent không
phân biệt rõ dữ liệu với instruction, nội dung bên ngoài có thể ảnh hưởng
đến quá trình reasoning và dẫn đến hành động sai.

BrowserART đại diện cho hướng đánh giá web agents trong môi trường trình
duyệt, nơi rủi ro có thể đến từ website, form, nội dung được render, hoặc
chỉ dẫn gián tiếp nằm trong trang web. Các benchmark dạng này giúp mở rộng
evaluation từ text-only jailbreak sang các tác vụ web thực tế hơn.

\subsection{So sánh các công trình tiêu biểu}
\label{subsec:tutorial-paper-comparison}

Bảng~\ref{tab:tutorial-paper-groups} tổng hợp các công trình tiêu biểu
được tutorial đề cập theo từng nhóm nội dung.

\begin{table}[H]
\centering
\caption{Các công trình tiêu biểu được tutorial đề cập theo từng nhóm nội dung.}
\label{tab:tutorial-paper-groups}
\begin{tabularx}{\textwidth}{
    >{\bfseries}p{0.23\textwidth}
    p{0.28\textwidth}
    X
}
\toprule
Nhóm nội dung
& Công trình tiêu biểu
& Vai trò trong tutorial \\
\midrule

White-box jailbreak
& Universal and Transferable Adversarial Attacks
& Đại diện cho adversarial suffix, thuật toán GCG và khả năng transfer giữa
các mô hình. \\

Black-box automated jailbreak
& PAIR; Tree of Attacks
& Đại diện cho attacker--target--judge, iterative refinement và search-based
jailbreak trong black-box setting. \\

Stealthy jailbreak generation
& AutoDAN
& Minh họa hướng tự động tạo jailbreak prompt tự nhiên hơn và khó phát hiện
hơn so với adversarial suffix bất thường. \\

Automated red teaming
& Red Teaming Language Models with Language Models; Adversarial Reasoning
at Scale
& Cung cấp nền tảng cho việc dùng mô hình ngôn ngữ hoặc test-time
computation để tự động tìm hành vi lỗi. \\

Inference-time defense
& SmoothLLM
& Đại diện cho defense không cần huấn luyện lại mô hình, dựa trên nhiều
perturbation của prompt và tổng hợp phản hồi. \\

Model-level defense
& Circuit Breakers; Representation Guardrails
& Đại diện cho hướng phát hiện hoặc can thiệp vào representation bên trong
mô hình để hạn chế hành vi không an toàn. \\

Evaluation benchmarks
& AdvBench; HarmBench; JailbreakBench; TDC
& Đại diện cho các tập hành vi, framework đánh giá automated red teaming,
robust refusal và benchmark mở cho jailbreak robustness. \\

Agentic systems
& AgentDojo; BrowserART; Breaking ReAct
& Đại diện cho đánh giá prompt injection, tool-use risks và security của
LLM agents trong môi trường có công cụ. \\

Robotics
& RoboPAIR
& Đại diện cho việc mở rộng jailbreak từ phản hồi văn bản sang hành động
trong robot được điều khiển bởi LLM. \\

\bottomrule
\end{tabularx}
\end{table}

Từ bảng trên, có thể thấy tutorial không chỉ trình bày jailbreak như một
tập prompt tấn công riêng lẻ. Thay vào đó, tutorial tổ chức lĩnh vực này
theo một chuỗi phát triển rộng hơn: từ attack thủ công sang attack tự động,
từ defense ở bề mặt input/output sang defense ở mức representation, từ
evaluation rời rạc sang benchmark chuẩn hóa, và từ chatbot sang agentic
systems hoặc robot.

\subsection{Nhận xét chung}
\label{subsec:tutorial-paper-discussion}

Từ các paper được tutorial đề cập, có thể rút ra bốn nhận xét chính.

Thứ nhất, jailbreak đang chuyển từ prompt engineering thủ công sang
automated red teaming. GCG, PAIR, Tree of Attacks, AutoDAN và các hướng
reasoning-based search cho thấy attacker có thể dùng tối ưu hóa, mô hình
ngôn ngữ hoặc search để tự động tìm prompt hiệu quả hơn.

Thứ hai, defense không thể chỉ dựa vào một lớp duy nhất. SmoothLLM cho thấy
có thể phòng thủ tại thời điểm suy luận; Circuit Breakers và Representation
Guardrails cho thấy khả năng phòng thủ ở tầng representation; nhưng mỗi
hướng đều có giới hạn về chi phí, quyền truy cập mô hình và khả năng chống
adaptive attacks.

Thứ ba, evaluation là yếu tố trung tâm để so sánh attack và defense.
AdvBench cung cấp tập harmful behaviors phổ biến; HarmBench chuẩn hóa
automated red teaming và robust refusal; JailbreakBench nhấn mạnh
artifacts, scoring functions và reproducibility. Điều này cho thấy một kết
quả ASR chỉ có ý nghĩa khi đi kèm threat model, target model, system prompt,
evaluator và ngân sách truy vấn.

Thứ tư, agentic systems làm thay đổi bản chất của bài toán. RoboPAIR,
AgentDojo, BrowserART và Breaking ReAct cho thấy khi LLM có memory, tools,
browser hoặc robot embodiment, rủi ro không chỉ nằm ở phản hồi cuối cùng.
Cần đánh giá toàn bộ trajectory, bao gồm observation, reasoning, tool call,
action và environment state.

\subsection{Tổng kết chương}
\label{subsec:tutorial-paper-summary}

Chương này đã tổng hợp các công trình tiêu biểu được đề cập trong tutorial
ICML 2025 về jailbreaking LLMs và agentic systems. Các paper về attack như
GCG, PAIR, Tree of Attacks và AutoDAN cho thấy jailbreak có thể được tối ưu
và tự động hóa. Các paper về defense như SmoothLLM, Circuit Breakers và
Representation Guardrails cho thấy phòng thủ có thể được triển khai ở nhiều
tầng khác nhau, từ runtime đến representation nội bộ của mô hình. Các
benchmark như AdvBench, HarmBench và JailbreakBench làm rõ nhu cầu chuẩn
hóa evaluation. Cuối cùng, RoboPAIR, AgentDojo, BrowserART và Breaking
ReAct mở rộng bài toán sang agentic systems, nơi mô hình có thể dùng công
cụ hoặc thực hiện hành động trong môi trường.

Những công trình này tạo nền tảng cho phần thực nghiệm ở
Chương~\ref{sec:experiments}, nơi báo cáo xây dựng một demo có kiểm soát
để minh họa attack--defense pipeline, cách đo lường safety--utility
trade-off và các giới hạn của đánh giá thực nghiệm.